% Standalone audit; compile directly with pdfLaTeX, twice.
% MAIN VERDICT: VERIFIED
% Exact P03 v1.0 source remains unchanged. No P04 mathematics is begun.
\documentclass[11pt]{article}
\usepackage[T1]{fontenc}
\usepackage[utf8]{inputenc}
\usepackage{lmodern,amsmath,amssymb,amsthm}
\usepackage[margin=0.85in]{geometry}
\usepackage{array,booktabs,longtable,enumitem,needspace}
\usepackage[hidelinks]{hyperref}
\newcolumntype{L}[1]{>{\raggedright\arraybackslash}p{#1}}
\setlength{\emergencystretch}{3em}
\setlength{\parindent}{0pt}
\setlength{\parskip}{5pt}
\setlist{nosep,leftmargin=*}
\newcommand{\RR}{\mathbb R}
\newcommand{\CC}{\mathbb C}
\newcommand{\ZZ}{\mathbb Z}
\newcommand{\one}{\mathbf 1}
\newcommand{\Dom}{(-1,\infty)^n}
\newcommand{\diag}{\operatorname{diag}}
\newcommand{\sa}{\operatorname{sa}}
\newcommand{\lin}[1]{\textbf{L#1}}
\newcommand{\sha}[2]{\par{\small\ttfamily #1\allowbreak #2}\par}
\title{A03: Independent mathematical audit of P03\\Origin and homogeneous linearization}
\author{Independent reconstruction and adversarial review}
\date{27 September 2026}
\begin{document}
\maketitle

\section{Main verdict and exact version}
\textbf{MAIN VERDICT: VERIFIED.}
Every substantive mathematical assertion in the exact P03 version identified
below passes independent reconstruction under its stated canonical assumptions.
There is no required mathematical or wording correction. The optional
nonzero-equilibrium simplification also passes within its stated domain;
it needs no separate verdict.

The audited file is \path{proofs/P03_origin_and_linearization.tex},
\textbf{P03 v1.0, 2026-09-27}, exactly \textbf{18,921 bytes and 433 lines},
from \path{cyclic_network_project_P03_draft.zip}. Its full-file SHA-256,
including comments, whitespace and final newline, is
\sha{0f87f315817429ddb949af5bc6d38db32}{a0cb6a8f7715dfe941840becd527776}
The supplied ZIP's SHA-256 is
\sha{768c76b7b53e0af2807d3eab81cbe381}{34425c3ff477585559428b9614d1fb96}
All P03 line references below use that exact source, including comment and
blank lines. The hash agrees with the manifest's final delivered-component
record, not its earlier pre-P03 input inventory. The supplied archive is the
version authority for this audit. No placeholder or archived blueprint is used
as the proof being audited.

Other local copies were checked for version confusion: the only substantive
P03 draft found has the same hash. Older 1,577-byte TODO placeholders have
hash beginning \texttt{8906056fbc57d1de} and were excluded.

The verdict rests on the analytic proofs in this document. Compilation,
finite samples and the earlier draft's validation report are not proof
premises. P03 and all supplied project files remain unchanged. A03 records
that this exact version passes; the unchanged P03 header still records its
pre-audit DRAFT state. Recording administrative audit closure is distinct from
altering the audited proof. P04 is cleared after that closure is recorded;
P04 is not begun here.

\section{Current inputs, assumptions and conventions}
All thirteen requested current inputs were located in the supplied archive
before reconstruction. The accompanying \path{input_hashes.json} records
their full SHA-256 digests and sizes.
\begin{longtable}{@{}L{0.51\textwidth}L{0.43\textwidth}@{}}
\toprule Current file & Use in this audit\\\midrule\endhead
\path{project_manifest.md} & Final P03 version/hash; P01/P02 closure; current status precedence.\\
\path{claim_registry.md} & C034, C035, C037 and the P03 overlay; distinction between archived claims and current statements.\\
\path{frozen_theorem_package.md} & S1 supporting role and the P03 then P04 attempt order.\\
\path{final_novelty_verdict.md} & Frozen literature scope and the separate, unclosed priority gate.\\
\path{02_theorem_dependency_map.tex} & P01/P02 to P03 to P04 dependency graph.\\
\path{00_canonical_model.tex} & Original vector field, parameter set, modular indexing and open domain.\\
\path{01_notation_and_assumptions.tex} & Predecessor shift, DFT sign, norms and hyperbolicity definition.\\
\path{proofs/P01_wellposedness_and_absorbing_box.tex} & Current v1.1 smoothness on the open domain.\\
\path{audits/A01_wellposedness_and_absorbing_box.tex} & Original audit and sole E01 correction, now closed.\\
\path{proofs/P02_synchronous_subspace_and_equilibria.tex} & Current v1.1 scalar restriction and nonzero-equilibrium relation.\\
\path{audits/A02_synchronous_subspace_and_equilibria.tex} & Restored original audit; sole E01 correction, now closed.\\
\path{proofs/P03_origin_and_linearization.tex} & Exact 433-line mathematical target.\\
\path{proofs_compendium.tex} & Current status banner, scope and unchanged module inclusion order.\\
\bottomrule
\end{longtable}

The field and the full standing assumptions are
\begin{align}
 F_i(x)&=a x_i-kx_i^2+\frac{c x_{i-1}}{1+x_i}
             -\frac{g x_i}{1+x_{i+q}}-h x_i x_{i-q},\label{A03:field}\\
 n&\in\ZZ,\quad n\ge5,\quad 1\le q\le n-1,\quad
 a,k,c,g,h>0,\quad x\in\Dom.\label{A03:assumptions}
\end{align}
Indices lie in $\ZZ/n\ZZ$ with representatives $0,\ldots,n-1$.
There is no coprimality assumption. Neighbor occurrences are retained even
when their coordinates coincide. In particular, the audit includes negative
coordinates in the open domain; it does not assume that this whole domain is
forward invariant. The present linear statements do not need such invariance.

Write $A=k+h$ and $\mu=a+c-g$. Exactly the project's conventions are used:
\[
 (Pz)_i=z_{i-1},\qquad (P^r z)_i=z_{i-r},\qquad
 \theta_m=\frac{2\pi m}{n},\qquad
 v^{(m)}_j=n^{-1/2}e^{\mathrm i\theta_m j}.
\]
The Fourier vectors have a positive exponent, so the predecessor shift has
a negative exponent in its eigenvalue. Hyperbolicity means that no eigenvalue
has zero \emph{real part}. All stability equivalences below concern the
autonomous linear equation $\dot z=J(0)z$ in finite-dimensional real space.

\section{Independent general Jacobian derivation}
Let $\delta_{j\ell}$ denote modular equality. Differentiate the five terms of
\eqref{A03:field} separately before merging any entries:
\begin{align*}
 \partial_j(a x_i)&=a\delta_{ji},\\
 \partial_j(-k x_i^2)&=-2kx_i\delta_{ji},\\
 \partial_j\!\left(\frac{c x_{i-1}}{1+x_i}\right)
 &=\frac{c}{1+x_i}\delta_{j,i-1}
    -\frac{c x_{i-1}}{(1+x_i)^2}\delta_{ji},\\
 \partial_j\!\left(-\frac{g x_i}{1+x_{i+q}}\right)
 &=-\frac{g}{1+x_{i+q}}\delta_{ji}
    +\frac{g x_i}{(1+x_{i+q})^2}\delta_{j,i+q},\\
 \partial_j(-h x_i x_{i-q})
 &=-h x_{i-q}\delta_{ji}-h x_i\delta_{j,i-q}.
\end{align*}
The denominator derivative of the positive input has an algebraic minus
sign; the denominator derivative of the negative loss has an algebraic plus
sign. Their numerical signs need not equal their algebraic signs when the
numerator coordinate is negative. All denominators here are nonzero.

Summing gives precisely
\begin{align}
 J_{ij}(x)={}&\left(a-2kx_i-\frac{c x_{i-1}}{(1+x_i)^2}
                  -\frac{g}{1+x_{i+q}}-h x_{i-q}\right)\delta_{ji}\notag\\
 &+\frac{c}{1+x_i}\delta_{j,i-1}
  +\frac{g x_i}{(1+x_{i+q})^2}\delta_{j,i+q}
  -h x_i\delta_{j,i-q}.\label{A03:J}
\end{align}
This verifies P03 \lin{33--70}. A modular Kronecker symbol is an indicator
appearing in a sum, not an instruction to overwrite an entry. Differentiation
with respect to a repeated coordinate differentiates every occurrence of it.
No application of a distinct-neighbor assumption is hidden in the derivation.

For an independent directional reconstruction, at fixed $x$ and $z$ use
$(u+t v)^{-1}=u^{-1}-t v/u^2+O(t^2)$ with $u>0$.
The rational first-order terms are
\[
 c\left(\frac{z_{i-1}}{1+x_i}
        -\frac{x_{i-1}z_i}{(1+x_i)^2}\right),\qquad
 -\frac{g z_i}{1+x_{i+q}}+\frac{g x_i z_{i+q}}{(1+x_{i+q})^2}.
\]
The other terms contribute $(a-2kx_i)z_i$ and
$-h(x_{i-q}z_i+x_i z_{i-q})$. This independently reproduces every coefficient
of \eqref{A03:J}, including at colliding offsets.

\section{Shift orientation and exhaustive local collisions}
Define $D_i,C_i,G_i,H_i$ by the four coefficients in \eqref{A03:J}, with
$H_i=h x_i$ so its contribution is $-H_i$. The directional derivative is
\[
 (DF(x)z)_i=D_i z_i+C_i z_{i-1}+G_i z_{i+q}-H_i z_{i-q}.
\]
From the indices themselves, $z_{i-1}=(Pz)_i$,
$z_{i+q}=(P^{-q}z)_i$, and $z_{i-q}=(P^qz)_i$. Therefore
\begin{equation}
 J(x)=\diag(D_i)+\diag(C_i)P+\diag(G_i)P^{-q}
                         -\diag(H_i)P^q.\label{A03:shift}
\end{equation}
The diagonal matrices multiply on the left because the coefficient belongs
to receiving row $i$. Right multiplication would generally sample a different
row's coefficient. This verifies P03 \lin{73--92} without deducing orientation
from its claimed answer. A nonsynchronous Jacobian need not be circulant.

No neighbor can equal $i$: $1\not\equiv0$ and $q\not\equiv0$ modulo $n$.
There are exactly three unordered pairs of neighbor occurrences to examine:
\begin{align*}
 i-1\equiv i-q&\iff q=1,\\
 i-1\equiv i+q&\iff q=n-1,\\
 i+q\equiv i-q&\iff n\mid 2q
                    \iff n\text{ even and }q=n/2.
\end{align*}
For the last equivalence, $0<2q<2n$ forces $2q=n$. The three values
$1,n-1,n/2$ cannot coincide for $n\ge5$; any triple collision would imply
two such cases and is impossible. These pairwise comparisons exhaust the
list. A nontrivial $\gcd(n,q)$ affects repeated step-$q$ paths, not these
one-row congruences, so it creates no additional local collision.

The actual general-state merged entries are
\begin{align*}
 q=1,\ j=i-1:&\quad \frac{c}{1+x_i}-h x_i,\\
 q=n-1,\ j=i-1:&\quad
              \frac{c}{1+x_i}+\frac{g x_i}{(1+x_{i-1})^2},\\
 n\text{ even},\ q=n/2,\ j=i+q=i-q:&\quad
              \frac{g x_i}{(1+x_{i+q})^2}-h x_i.
\end{align*}
In each case the remaining neighbor contribution is retained. A merged entry
may cancel to zero. This does not change the differentiation rule.

\section{Origin derivative, complete spectrum and multiplicities}
All five terms vanish at $x=0$, so $F(0)=0$. In the derivative:
\begin{itemize}
\item $-2kx_i$ and both $h$-term derivatives vanish because they contain a state coordinate.
\item The $c$ denominator derivative vanishes because it contains $x_{i-1}$.
\item The $g$ denominator derivative vanishes because it contains $x_i$.
\item The numerator derivatives survive as $c$ at $i-1$ and $-g$ at $i$; the linear term contributes $a$ at $i$.
\end{itemize}
Hence
\begin{equation}
 J(0)=(a-g)I+cP.\label{A03:origin}
\end{equation}
This includes every collision: for $q=1$ the merged coefficient is $c-0$;
for $q=n-1$ it is $c+0$; for $q=n/2$ the half-offset coefficient is $0-0$.
No diagonal collision can revive a vanished term.

Apply $P$ directly to the prescribed Fourier vector:
\[
 (Pv^{(m)})_j=n^{-1/2}e^{\mathrm i\theta_m(j-1)}
             =e^{-\mathrm i\theta_m}v^{(m)}_j.
\]
Wraparound is valid since $e^{\mathrm i n\theta_m}=1$. The geometric-series
identity shows $(v^{(\ell)})^*v^{(m)}=\delta_{\ell m}$, so the $n$ vectors
are an orthonormal basis. Equation \eqref{A03:origin} gives the full spectrum
\begin{align}
 \lambda_m&=a-g+c e^{-\mathrm i\theta_m},\quad 0\le m<n,\label{A03:spectrum}\\
 \Re\lambda_m&=a-g+c\cos\theta_m
              =\mu-c(1-\cos\theta_m),\notag\\
 \Im\lambda_m&=-c\sin\theta_m.\notag
\end{align}
Distinct roots of unity remain distinct after multiplying by $c>0$ and
translating. Consequently each of the $n$ eigenvalues has algebraic and
geometric multiplicity one over $\CC$. Equal real parts of conjugate modes
are not repeated eigenvalues. This argument concerns the origin only.

An independent determinant route avoids assuming the spectral formula.
Write $t=\lambda-(a-g)$. In the determinant expansion of $tI-cP$, each row
can choose its diagonal or predecessor column. If any row chooses its
predecessor column, column injectivity forces the preceding row to choose
its predecessor too; going around the cycle forces all rows to do so.
Thus the only contributing permutations are the identity and the full
$n$-cycle. Their contributions are $t^n$ and
$(-1)^{n-1}(-c)^n=-c^n$, respectively. Therefore
\begin{equation}
 \chi(\lambda)=(\lambda-(a-g))^n-c^n.\label{A03:characteristic}
\end{equation}
At a root, $|\lambda-(a-g)|=c>0$, so
$\chi'(\lambda)=n(\lambda-(a-g))^{n-1}\ne0$. This independently confirms
algebraic simplicity; the already constructed eigenvector gives geometric
multiplicity one. P03 \lin{94--158} passes.

\section{Spectral abscissa and exact linear stability}
For $0\le m<n$, $\cos\theta_m\le1$, with equality exactly when $m=0$.
Because $c>0$,
\[
 \sa(J(0))=\mu,\qquad
 \lambda_0-\Re\lambda_m=c(1-\cos\theta_m)>0\quad(m\ne0).
\]
The maximizing mode is unique in every odd or even canonical dimension.
Neither $q$ nor its gcd with $n$ enters \eqref{A03:origin}; the spectral
abscissa and its sign do not even depend on $n$. Individual nonzero-mode
eigenvalues and their zero-real-part loci do depend on $n$.

Let $V$ have columns $v^{(m)}$ and put $\widehat z=V^*z$. Orthonormality and
the origin eigenvectors give $\widehat z_m(t)=e^{t\lambda_m}\widehat z_m(0)$.
For real initial data, conjugate Fourier coefficients preserve reality.
Thus for $t\ge0$,
\[
 \|z(t)\|_2^2=\sum_m e^{2t\Re\lambda_m}|\widehat z_m(0)|^2
             \le e^{2\mu t}\|z(0)\|_2^2.
\]
If $\mu<0$, this proves exponential stability with constant $1$ and decay
rate $-\mu$. If $\mu>0$, arbitrarily small real data $\varepsilon\one$
produce $\varepsilon e^{\mu t}\one$ and leave any fixed sufficiently small
ball, proving Lyapunov instability.

If $\mu=0$, the norm is nonincreasing, every nonuniform Fourier component
decays, and the uniform component is constant. This proves Lyapunov
stability but rules out attraction, asymptotic stability and exponential
stability. The three cases exhaust the real parameter $\mu$, proving both
directions of all P03 equivalences. In finite dimension these properties
do not depend on the choice of norm. No nonlinear stability theorem is used
or concluded, especially at $\mu=0$.

\section{Full hyperbolicity-loss set and imaginary pairs}
By definition and \eqref{A03:spectrum}, the exact nonhyperbolicity set in the
canonical parameter domain is
\begin{equation}
 \bigcup_{m=0}^{n-1}
 \left\{(a,k,c,g,h)>0:
       g=a+c\cos\frac{2\pi m}{n}\right\}.\label{A03:nonhyp}
\end{equation}
Only positive values of the right-hand side are admissible. A proposed locus
with $a+c\cos\theta_m\le0$ contributes no canonical point at those $a,c$.
Modes $m$ and $n-m$ name the same locus, not two different parameter
conditions. Hyperbolicity holds exactly on the complement of this union.

At every nonuniform locus,
$\mu=c(1-\cos\theta_m)>0$. Thus the origin can be both nonhyperbolic and
linearly unstable. The stability boundary is only the $m=0$ locus;
the other components must not be discarded.

For an angular interior representative $0<m<n/2$, one has
$0<\theta_m<\pi$, hence $\sin\theta_m>0$. On its admissible locus,
\[
 \lambda_m=-\mathrm i c\sin\theta_m\ne0,\qquad
 \lambda_{n-m}=+\mathrm i c\sin\theta_m\ne0.
\]
These are distinct simple eigenvalues. On $[0,\pi]$ cosine is strictly
decreasing. Choosing representatives $0\le j\le\lfloor n/2\rfloor$
therefore allows only one representative at a prescribed zero-real-part
value, and exactly the two modes $m,n-m$ are critical for this interior case.
All other modes have nonzero real part; this does \emph{not} assert they
are all stable. In fact mode zero is already unstable.

For odd $n$, the highest representative $(n-1)/2$ still satisfies
$0<m<n/2$ and belongs to a pair. For even $n$, $m=n/2$ is a singleton
with zero sine. There is no coincidence of two distinct Fourier pairs at
one such locus for fixed $a,c>0$. This settles both parities. P03
\lin{168--206} passes. These facts are not a Hopf bifurcation theorem.

\section{Zero eigenvalues: boundary and nonuniform singleton}
The condition $\lambda_m=0$ is stronger than $\Re\lambda_m=0$. Since
$c>0$, its imaginary part forces $\sin\theta_m=0$. Among canonical mode
labels this means $m=0$, or, only for even $n$, $m=n/2$.

At the stability boundary $g=a+c$,
\[
 J(0)=c(P-I),\qquad
 \lambda_m=c(e^{-\mathrm i\theta_m}-1).
\]
Only $m=0$ has zero real part, and only that mode has eigenvalue zero.
Every other real part equals $-c(1-\cos\theta_m)<0$. The equation
$Pz=z$ forces all coordinates equal, so the kernel is
$\operatorname{span}\{\one\}$ (over $\RR$ or $\CC$, respectively).
Independently, $\chi(\lambda)=(\lambda+c)^n-c^n$ has
$\chi'(0)=n c^{n-1}\ne0$. Algebraic and geometric multiplicities are both
one. For even $n$, the singleton eigenvalue is $-2c$, not zero.

For a nonzero-mode zero, the only remaining possibility is even $n$ and
$m=n/2$, where $\lambda_{n/2}=a-g-c$. Hence its necessary and sufficient
condition is
\[
 n\text{ even},\quad m=n/2,\quad g=a-c>0,
 \quad\text{equivalently }a>c\text{ and }g=a-c.
\]
Here $\mu=2c>0$, so the synchronous mode is unstable. Also $J(0)=c(I+P)$;
$Pz=-z$ gives the one-dimensional alternating kernel
$\operatorname{span}\{((-1)^j)_{j=0}^{n-1}\}$. Even parity is necessary
for this recurrence to close. The characteristic polynomial is
$(\lambda-c)^n-c^n$ and its derivative at zero is
$n(-c)^{n-1}=-n c^{n-1}\ne0$. Thus this zero is algebraically and
geometrically simple as well. Odd dimensions have no nonuniform zero.
P03 \lin{208--250} passes without extending its nonlinear scope.

\section{Homogeneous coefficients and equilibrium-only simplification}
Let $x=r\one$ with $r>-1$ and set $s=1+r>0$. Substitute into the
independently derived \eqref{A03:J}, without assuming an equilibrium:
\begin{align}
 d&=a-2kr-\frac{cr}{s^2}-\frac{g}{s}-hr,\notag\\
 C&=\frac{c}{s},\qquad G=\frac{gr}{s^2},\qquad H=hr,\label{A03:homcoef}\\
 J(r\one)&=dI+CP+GP^{-q}-HP^q.\label{A03:homJ}
\end{align}
The directional action is $d z_i+Cz_{i-1}+Gz_{i+q}-Hz_{i-q}$, which
again checks the signs and orientations directly. These formulas hold at
arbitrary synchronous states, including negative ones.

Always $C>0$. For $r>0$, $G,H>0$ and the separate neighbor signs are
$+,+,-$. For $r=0$, $G=H=0$. For $-1<r<0$, $G,H<0$, so the separate
neighbor signs are $+,-,+$. At an arbitrary synchronous state there is
no universal sign of $d$; for example at $r=0$ it is $a-g$. The sign of a
merged neighbor entry is decided only after summation.

Only if $r\ne0$ and $r\one$ is an equilibrium does P02 permit division by
$r$ to yield $a-(k+h)r+(c-g)/s=0$. Substituting
$a=(k+h)r-(c-g)/s$ into \eqref{A03:homcoef} gives
\begin{align*}
 d&=(k+h)r-\frac{c-g}{s}-2kr-\frac{cr}{s^2}-\frac{g}{s}-hr\\
  &=-kr-\frac{c}{s}-\frac{cr}{s^2}
   =-kr-\frac{c(1+2r)}{s^2}.
\end{align*}
The $h$ and $g$ cancellations are exact. This works for positive or
negative nonzero equilibria. It is not valid at arbitrary synchrony.
At the origin the unsimplified value is $a-g$, whereas the displayed
simplified expression would be $-c$; they agree only if the additional
condition $g=a+c$ is separately imposed. Merely being the origin equilibrium
does not supply that condition. P03 makes exactly this restriction at
\lin{295--311}, so no correction is needed.

\section{Homogeneous collisions and circulant structure}
Use the exhaustive congruence list from Section 4 in \eqref{A03:homJ}:
\begin{align*}
 q=1:&\qquad J=dI+(C-H)P+GP^{-1},\\
 q=n-1:&\qquad J=dI+(C+G)P-HP^{-1},\\
 n\text{ even},\ q=n/2:&\qquad J=dI+CP+(G-H)P^{n/2}.
\end{align*}
In the second case $P^{-q}=P$ and $P^q=P^{-1}$; in the third,
$P^{-q}=P^q=P^{n/2}$. All signs follow from the original directional
derivative. The cases are mutually exclusive for $n\ge5$ and no triple
collision occurs. Zero or cancelling coefficients cause no exception.

Since $P^n=I$, the general homogeneous expression is the ordinary polynomial
$dI+CP+GP^{n-q}-HP^q$. For every integer $r$, the entry of $P^r$ in row
$i$, column $j$ is one exactly when $i-j\equiv r$. Thus every entry of the
polynomial depends only on $i-j$ modulo $n$, proving the circulant property
at every synchronous state. Linear combination remains valid if powers
coincide or their coefficients cancel.

P03's mention of the standard Fourier-diagonalizability consequence does
not supply or use the general homogeneous eigenvalue formula. Its full
homogeneous spectrum and normality theorem are expressly deferred to P04.
Neither is derived or promoted here. P03 \lin{313--355} passes.

\section{Independent synchronous-direction consistency check}
Starting directly from P02's scalar field, with a fresh state variable $r$,
\[
 f(r)=ar-(k+h)r^2+\frac{(c-g)r}{1+r},\qquad
 f'(r)=a-2(k+h)r+\frac{c-g}{(1+r)^2}.
\]
For $s=1+r$, sum the unsimplified coefficients:
\begin{align*}
 d+C+G-H
 &=a-2(k+h)r+\frac{c-g}{s}+\frac{(g-c)r}{s^2}\\
 &=a-2(k+h)r+\frac{(c-g)(s-r)}{s^2}=f'(r).
\end{align*}
The identity $s-r=1$ completes the check. Equivalently differentiate
$F(r\one)=f(r)\one$ with respect to $r$: $J(r\one)\one=f'(r)\one$.
Every shift fixes $\one$, so collisions preserve the same sum. No equilibrium
condition is used. This is a directional/row-sum check, not a derivation of
the full positive-equilibrium spectrum. P03 \lin{375--385} passes.

\section{Adversarial exact cases and counterexample search}
The universal arguments above establish the verdict. The following exact
cases actively test likely failure mechanisms; they are not substitutes for
those arguments. Unspecified $k,h$ in origin examples can both be set to $1$.

\subsection*{Origin, parity and hyperbolicity stress cases}
Take $a=2,c=1$. At $g=3$, every $n=5,6,8$ has one simple zero at $m=0$
and all other real parts strictly negative. At $g=4$ all modes decay; at
$g=2$ the synchronous growth rate is $1$. At $g=1$, every even dimension
has its simple alternating zero with $\mu=2$, whereas $n=5$ has no zero
eigenvalue. These examples apply to every canonical $q$, including collisions.

For interior loci with the same $a,c$, the following table gives exact
critical values. Each listed mode is paired with $n-m$; only that pair has
zero real part. Set $\omega=-\Im\lambda_m>0$.
\begin{center}
\begin{tabular}{ccccc}
\toprule $n$ & $m$ & $g$ & $\mu$ & $\omega$\\\midrule
5 & 1 & $2+(\sqrt5-1)/4$ & $(5-\sqrt5)/4$ & $\sqrt{10+2\sqrt5}/4$\\
5 & 2 & $2-(\sqrt5+1)/4$ & $(5+\sqrt5)/4$ & $\sqrt{10-2\sqrt5}/4$\\
6 & 1 & $5/2$ & $1/2$ & $\sqrt3/2$\\
6 & 2 & $3/2$ & $3/2$ & $\sqrt3/2$\\
8 & 1 & $2+\sqrt2/2$ & $1-\sqrt2/2$ & $\sqrt2/2$\\
8 & 2 & $2$ & $1$ & $1$\\
8 & 3 & $2-\sqrt2/2$ & $1+\sqrt2/2$ & $\sqrt2/2$\\
\bottomrule
\end{tabular}
\end{center}
All these $g$ values are positive. The $n=5,m=2$ row specifically checks
the largest odd-dimensional representative: it is a nonreal pair, not an
even-style singleton. Equal frequencies in different rows do not create
simultaneous criticality, since their $g$ values differ.

To test the positivity filter, set $a=1/4,c=1,n=6$. The formal $m=2$
locus has $g=-1/4$ and is inadmissible; the alternating locus has $g=-3/4$
and is also inadmissible. They cannot be added to the canonical set. Conversely
the $m=1$ locus has $g=3/4>0$ and is retained.

\subsection*{Collisions, gcd, signed states and cancellation}
For $n=5$, $q=1,4$ test the two predecessor collisions, while $q=2,3$
give distinct neighbors. For $n=6$, $q=3$ gives the half-offset collision;
$q=2,4$ have nontrivial gcd but distinct neighbor offsets. For $n=8$,
$q=4$ gives the half-offset collision, whereas $q=2,6$ have gcd $2$
without a collision. Taking $q=1,n-1$ in these even dimensions tests the
other two cases and their exclusivity.

Two positive nonzero equilibria deliberately force merged coefficients to
vanish. At $r=1$ and $(a,k,c,g,h)=(5/2,1,2,3,1)$,
\[
 d=-5/2,\quad C=H=1,\quad G=3/4.
\]
For $(n,q)=(5,1)$, $J=-\tfrac52 I+\tfrac34P^{-1}$: the predecessor
contribution cancels completely. At $r=1$ and
$(a,k,c,g,h)=(7/2,1,1,4,1)$,
\[
 d=-7/4,\quad C=1/2,\quad G=H=1.
\]
For $(n,q)=(6,3)$ or $(8,4)$, $J=-\tfrac74 I+\tfrac12P$.
Each tuple satisfies the equilibrium relation before it is used to simplify
$d$. These examples refute any unstated requirement that every listed
cyclic diagonal be nonzero, but they do not refute P03.

Cancellation of $C+G$ is possible at negative synchrony as well: take
$(n,q)=(6,5)$, $r=-1/2$ and $(a,k,c,g,h)=(4,1/2,2,2,1)$.
Then $(d,C,G,H)=(5,4,-4,-1/2)$ and $J=5I+\tfrac12P^{-1}$.
This is an arbitrary synchronous state, not an asserted equilibrium.
P03's additive formulas cover it; its two named cancellation types are
examples, not a purported exhaustive classification of cancellations.

A negative nonzero equilibrium in the open domain is $r=-1/2$ with
$(a,k,c,g,h)=(1,1,1,2,1)$. It has
\[
 d=1/2,\quad C=2,\quad G=-4,\quad H=-1/2,\quad f'(r)=-1.
\]
Indeed $a-(k+h)r+(c-g)/(1+r)=1+1-2=0$ and
$d+C+G-H=-1$. The three possible merged coefficients are respectively
$C-H=5/2$, $C+G=-2$, and $G-H=-7/2$. This tests negative-state sign
reversals and confirms that $d$ need not be negative even at a negative
equilibrium. For comparison, the same $r$ with all five parameters equal
to $1$ is not an equilibrium: its unsimplified $d=5/2$, while the
equilibrium-only formula would give $1/2$. At the origin with all five
parameters equal to $1$, the correct $d=0$, while that formula would give
$-1$. These are adversarial controls against misuse, not errors in P03.

\subsection*{Reproducible finite corroboration}
The accompanying exact checker differentiates the original five-term field
using rational automatic differentiation, with actual modular indices
identified first. It compares the result against the independently written
entry and shift expressions for all $q$ in dimensions $5,6,8$, including
nonsynchronous signed rational states, positive and negative synchrony,
the origin, nonzero equilibria and cancellations. All arithmetic in that
checker is exact; its script and machine-readable results are included in
the audit-only package. The results report actual counts and checks.
The hand proofs remain necessary for arbitrary $n,q,x$ and parameters.

The independent run passes 656 full Jacobian cases, comprising 29,820 exact
entry comparisons over all 16 dimension/offset pairs. It includes 456
homogeneous cases (20,734 matrix entries), 3,020 scalar row-sum comparisons,
48 origin cases, and 112 actual nonzero equilibria (64 positive and 48
negative). At those equilibria 742 original field components vanish and
742 diagonal simplifications agree exactly. All three merged-coefficient
cancellation types are covered. Counts for overlapping identities are not
added as if they were independent samples. Separate exact sparse-determinant
and rational-rank checks at $n=5,6,8$ corroborate the origin calculation.

The earlier P03 computational paragraph is explicitly corroborative and
is not imported as evidence of proof. Its claimed sample counts are
arithmetically consistent:
\[
 \sum_{n=5}^{10}(n-1)=39,\quad
 \sum_{n=5}^{10}(n-1)n^2=2570,\quad
 8\cdot39=312,\quad8\cdot2570=20560.
\]
Historical residuals are properties of their recorded program and environment,
not universal mathematical error bounds. The independent checking record
distinguishes those reports from the exact checks performed for A03.

The matching original companion was located at
\path{output/p03/P03_validation_support.zip}. Its checker has the manifest's
SHA-256
\sha{73784fd269f445a80639175b8901610e}{89a517469d69ddc5d397200dc5412bbb}
The inspected checker was copied and rerun without modifying the original
reports or reference files. In Python 3.12.14 and NumPy 2.3.5, its exact and
floating result records match the originals, including the maximum scaled
finite-difference error $2.6907607135803086\times10^{-9}$, correctly rounded
up to $2.70\times10^{-9}$ in P03. The maximum absolute difference is
$7.230752174791633\times10^{-8}$; the origin eigenvector residual is
$1.1889678785578623\times10^{-15}$. The reproduction comparison and summary
are retained with A03. This confirms the reported finite experiment, not
the universally quantified theorem statements.

No counterexample to a mathematical P03 claim was found. More importantly,
the derivations above explain why the attempted failures cannot occur under
the canonical hypotheses.

\section{Terminology, scope and dependency verification}
P03 uses ``asymptotically'', ``exponentially'', ``Lyapunov stable'' and
``unstable'' in \lin{163--193} for the explicitly specified linear system.
It does not claim that nonlinear asymptotic stability is equivalent to
$\mu<0$. At $\mu=0$, it classifies only the linear system. ``Nonhyperbolic''
in \lin{168--176,244--250} is correctly based on zero real part, and is
allowed to overlap linear instability. The interior imaginary eigenvalues
in \lin{198--206} are nonzero. The zero-eigenvalue results in
\lin{208--242} impose the additional zero-imaginary-part requirement.
No conflation is present.

The only dependencies used are the canonical field and conventions, P01's
smoothness on the open domain, P02's scalar restriction and explicitly invoked
nonzero-equilibrium relation, elementary differentiation, elementary cyclic
permutation/Fourier facts, and scalar linear differential equations.
P01's uniqueness supports P02's synchronous restriction but is not an extra
spectral premise. P01's global absorbing-box conclusions and P02's full root
classification are not needed by P03's linear algebra.

Current P01 v1.1 has full-file hash
\sha{0d6950821aa83f3d09e8bfe84dab2b0b}{5020aa6e7881cf6928e984bff1890f1e}
and current P02 v1.1 has full-file hash
\sha{8e057c2cd009659eee45906844e25f4f}{38c156e5eb2e4130e824aa5e8299c51a}
Their current corrections are present: P01's larger-real-root wording and
P02's explicitly ordered $(s_-,s_+)$ boundary limits. A01/A02 retain their
original verdicts on earlier exact versions; the later closure records,
current headers and manifest establish the accepted prerequisite state.
Restored A02 source is present despite P02's preserved historical availability
comment. Historical all-TODO and DRAFT sentences are not treated as current
status instructions.

P03 does not use a P04 theorem to prove its origin result. Its origin Fourier
calculation is supplied in full. The old notation document's broad ownership
sentence about conjugacy/simplicity must be read with the current P03 origin
scope, which is explicit in the current manifest, registry and compendium.
This is historical scope bookkeeping, not a circular mathematical dependency.

No full positive-equilibrium Fourier spectrum, homogeneous normality theorem,
positive-state mode selection, nonlinear boundary stability, transcritical
bifurcation, Hopf bifurcation, first Lyapunov coefficient or rotating wave is
proved or claimed complete in P03. Origin spectral ordering is part of its
authorized scope. P03's explicit deferrals at \lin{22--25,244--250,320--322}
are respected. No new literature search was conducted, and the frozen
priority limitation remains unchanged.

\section{Line-by-line issue and disposition list}
\textbf{Required issue list: empty.} The following contiguous line ranges
cover every line of the exact 433-line source, including nonmathematical
headers, comments and whitespace. ``Pass'' is a mathematical/scope finding;
``Record'' identifies status or computational provenance rather than a theorem.
\small
\begin{longtable}{@{}L{0.12\textwidth}L{0.12\textwidth}L{0.70\textwidth}@{}}
\toprule P03 lines & Disposition & Independent finding\\\midrule\endhead
1--11 & Record & Correct v1.0 DRAFT provenance, prerequisite and origin-only scope labels. No numerical dependency asserted.\\
12--26 & Pass & Canonical assumptions, domain and abbreviations exact; linear/nonlinear distinction and P04/P09 deferrals explicit.\\
27--32 & Pass & Modular indicators correctly defined as additive.\\
33--48 & Pass & General entry formula correct; neighbors off diagonal but potentially mutually coincident.\\
49--59 & Pass & Both rational derivatives correct, including opposite denominator-derivative signs.\\
60--71 & Pass & Loss-denominator derivative has sign of $x_i$; polynomial derivatives and additive handling correct.\\
72--79 & Pass & $D_i,C_i,G_i,H_i$ match the separately differentiated terms.\\
80--90 & Pass & Directional formula and $P,P^{-q},P^q$ orientations exact.\\
91--93 & Pass & Left diagonal factors use receiving-row coefficients; no general-state circulant claim.\\
94--104 & Pass & Origin equilibrium and every vanishing/surviving derivative identified correctly.\\
105--109 & Pass & $J(0)=(a-g)I+cP$, including collisions.\\
110--118 & Pass & Complete origin spectrum with negative exponent; all eigenvalues simple over $\CC$.\\
119--130 & Pass & Real/imaginary parts, unique maximizing mode and spectral abscissa correct.\\
131--142 & Pass & Componentwise shift, wraparound and geometric-sum basis proof correct.\\
143--152 & Pass & Distinctness, characteristic polynomial and nonzero derivative prove multiplicities.\\
153--159 & Pass & Strict cosine inequality for all nonzero modes, all parities.\\
160--177 & Pass & Linear stability trichotomy and full nonhyperbolicity set correct; $g>0$ retained.\\
178--193 & Pass & Fourier solution, norm bounds, real growing solution and both converses correct.\\
194--200 & Pass & Independence from $q$/gcd and threshold independence from $n$ correct; $n=6$ example exact and no Hopf assertion.\\
201--207 & Pass & Exactly one angular interior conjugate pair can be critical. For odd $n$, the largest representative is still interior in $0<m<n/2$.\\
208--214 & Pass & Boundary zero unique, simple, with synchronous kernel; other modes strictly left.\\
215--229 & Pass & Direct kernel, determinant derivative and even singleton $-2c$ establish every claim.\\
230--242 & Pass & Nonuniform zero iff even $n$, $m=n/2$, $g=a-c>0$; simple, with synchronous instability.\\
243--251 & Pass & Nonhyperbolic/unstable overlap and linear/nonlinear distinction correctly stated; no transcritical claim.\\
252--257 & Pass & All synchronous states $x^*>-1$ included; local use of $s=1+x^*$ disclosed.\\
258--276 & Pass & Unsimplified homogeneous coefficients and summed entries exact without equilibrium assumptions.\\
277--288 & Pass & Direct substitution and index-based shift proof valid.\\
289--294 & Pass & Positive, zero and negative-state sign qualifications correct; signs after collisions not assumed.\\
295--301 & Pass & P02 relation used only for a nonzero synchronous equilibrium.\\
302--312 & Pass & Both simplified forms of $d$ exact; arbitrary states and automatic use at origin explicitly excluded.\\
313--323 & Pass & Polynomial in $P$ proves circulant structure; generic spectrum and normality deferred.\\
324--335 & Pass & Pairwise congruence list exhaustive; diagonal/triple collisions excluded; no extra gcd case.\\
336--348 & Pass & Three homogeneous collision forms, exclusivity and possible zero coefficients correct.\\
349--356 & Pass & General-state merged entries use the identified coordinate and add all terms; cancellation allowed.\\
357--374 & Pass & First-order inverse expansion and independent origin substitution reproduce derivative.\\
375--386 & Pass & Scalar derivative equals row sum without the equilibrium relation.\\
387--393 & Pass & Independent multiplicity/order checks valid; equality of real parts distinguished from eigenvalue multiplicity.\\
394--409 & Record & Finite corroboration explicitly not proof. Counts are consistent; reproduced checks and provenance are recorded separately. No later spectrum or nonlinear result is inferred.\\
410--432 & Record & Draft auditor checklist is reference material. All relevant requested targets are addressed by this independently commissioned audit.\\
433 & Record & Historical DRAFT-status statement remains unchanged; A03 supplies the new independent verdict.\\
\bottomrule
\end{longtable}
\normalsize

One optional clarity note, not a required revision: at \lin{201--206},
``interior representative'' is most explicitly read as $0<m<n/2$.
For odd $n$, this includes the largest representative $\lfloor n/2\rfloor$.
The full all-mode formula and the project's singleton convention already
give the correct result. This reading is verified explicitly above and does
not change the verdict.

\section{Downstream consequences and closure}
The exact P03 v1.0 identified in Section 1 has passed the independent
mathematical audit. There is no unresolved correction preventing closure.
The accepted outputs are the general Jacobian with additive collisions,
the complete origin spectrum and linear classification, the full origin
nonhyperbolicity/zero-eigenvalue analysis, and the unsimplified homogeneous
circulant representation with its correctly restricted simplification.

\textbf{P04 is cleared to begin after A03 closure is recorded against this
exact P03 hash.} Closure can cite this audit and its VERIFIED verdict while
preserving the mathematical P03 source. The manifest/registry/compendium
have not been changed in this audit-only delivery. If a future administrative
header change alters P03's full-file hash, that closure record should preserve
the audited hash and explicitly account for the metadata-only difference.
A substantive proof change requires review of the changed version.

Clearance is permission to attempt P04, not verification of P04 or of a
positive-equilibrium spectrum, normality, bifurcation or later-module claim.
P09 still has its other dependencies; this audit does not decide nonlinear
behavior at $\mu=0$. The frozen literature-priority limitation remains a
separate publication gate. No P04 work or new literature work is performed.

The audit-only package contains the TeX source, PDF reading copy, input
fingerprints, reproducible checks and provenance. Every supplied project
member is preserved byte-for-byte.
\end{document}
